\section{Conclusions}
\label{sec:conclusion}

This paper presents a reproducible selective-adaptation workflow for scarce-label remote sensing scene classification. Built on DINOv2-Base, the workflow combines selective finetuning, residual adapters, feature-center regularization, fixed split manifests, and scripted multi-seed reporting. Under the evaluated 20-shot protocol on AID and NWPU-RESISC45, it improves over a matched-split ConvNeXt-Tiny reference in three-seed mean overall accuracy while remaining trainable on a single 8 GB GPU. The main contribution is therefore a rebuildable geocomputing recipe for public-data low-shot evaluation, not a new remote-sensing-specific architecture.

The current evidence suggests that selective backbone finetuning is the main source of mean improvement, whereas the residual adapter mainly improves stability and the feature-center term mainly improves mean performance under the current protocol. The resulting contribution is therefore not a claim of lightweight inference or universal superiority, but a practical geocomputing recipe that can be rebuilt, audited, and extended on public data and modest hardware.

Several limitations remain. The gain on NWPU-RESISC45 is positive but modest, so the cross-dataset conclusion should be interpreted conservatively and read mainly as evidence that the workflow remains stable under a second public benchmark. The method is parameter-efficient in trainable subset rather than lightweight in total inference complexity. Finally, the current study covers two public datasets under a fixed 20-shot protocol; broader shot regimes and additional remote sensing domains would strengthen future validation. Even with these limits, the workflow provides a useful open baseline for scarce-label geospatial image analysis and a concrete starting point for follow-up geocomputing studies.
